\documentclass[10pt,a4paper]{amsart}
\usepackage[margin=19mm]{geometry}
\usepackage{amsmath,amssymb,mathtools}
\usepackage[scr=rsfs,cal=euler]{mathalfa}
\usepackage{xfrac}
\usepackage[unicode,hidelinks,bookmarks=false]{hyperref}
\newcommand{\ie}{i.e.\ }
\newcommand{\cf}{cf.\ }
\newcommand{\resp}{respectively\ }
\newcommand{\Subsection}{\subsection}
\providecommand{\nobreakdash}{\nobreak\hbox{-}}
\setlength{\emergencystretch}{2em}
\multlinegap0pt
\usepackage{amssymb}
\newtheorem{lemma}{Lemma}
\title{The biharmonic hypersurface flow and the Willmore flow in higher dimensions}
\author{Yu Fu, Min-Chun Hong, Gang Tian}
\date{Source: arXiv:2505.19727v1 (CC BY 4.0)}
\begin{document}
\maketitle

\noindent\emph{Source and licence.} The biharmonic hypersurface flow and the Willmore flow in higher dimensions. Version arXiv:2505.19727v1 (CC BY 4.0), distributed by arXiv under the Creative Commons Attribution 4.0 International licence, http://creativecommons.org/licenses/by/4.0/, as recorded in the arXiv licence metadata for this exact version and shown on the abstract page https://arxiv.org/abs/2505.19727v1. Full licence notice: \texttt{licenses/arxiv-2505.19727-CC-BY-4.0.txt}.\par\smallskip

\noindent\textit{Editorial scope.} Counted unit: the article's lemma L1 (source0.tex lines 567--578). The article's complete original proof of it is delivered with it (source0.tex lines 580--695, one \texttt{proof} environment of 116 lines). No mathematics is written, repaired or re-derived: the statement and the proof are the article's own lines, in the article's own notation, and no retained line is substituted or re-flowed.  Retained uncounted as the source-local context the proof needs: lines 394--398.  Nothing was omitted from any retained span, so the package carries no omission record.  No retained line contains a citation, so the wrapper carries no bibliography and no bibliography entry.  No retained line contains a figure environment, an image-inclusion command or drawing code, so no image is delivered and the package declares no asset dependency.  Editorial formatting only, in addition: an AMS \texttt{amsart} preamble that restates the article's own declarations and macros used by the retained lines, namely the environments \texttt{lemma} on separate counters, together with the standard packages those lines need.  The retained environments print in this document as Lemma 1 in order of appearance, whereas the article prints the same items with the numbering of its own full document.  The complete native source of the article as distributed in the arXiv e-print for this exact version is bundled unchanged as \texttt{sources/178767/source0.tex}.
\begin{align}\label{GW}
\frac{\partial^2 f}{\partial x_i \partial
    x_j}=\Gamma_{ij}^k\frac{\partial f}{\partial x_k}+A_{ij}\nu,\quad
\frac{\partial}{\partial x_i} \nu=-{{}\bf S} (\frac{\partial f}{\partial x_i}),
\end{align}

\begin{lemma} \label{L1}
Let $f: M^n\times [0, T] \rightarrow \mathbb{R}^{n+1}$ be a solution to the flow $ \frac{\partial f}{\partial t}=F \nu+W$ with general $F$ and $W= a^l\frac {\partial f}{\partial x_l }$. Then we have
\begin{align}
\label{L21}\frac{\partial}{\partial t} g_{ij} &= -2F A_{ij}+\nabla_i a^l g_{jl}+\nabla_j a^lg_{il} \mbox{,}\\
\label{L22}\frac{\partial}{\partial t} g^{ij} &= 2F A^{ij}-\nabla_l a^i g^{jl}-\nabla_l a^jg^{il} \mbox{,}\\
\label{L23}\frac{\partial}{\partial t} \nu &=-\nabla F-a^l g^{ij}A_{il}\partial f_j \mbox{,}\\
\label{L24}\frac{\partial}{\partial t} d \mu &= (-HF+\nabla_l a^l)~d\mu  \mbox{,}\\
\label{L25}\frac{\partial}{\partial t} A_{ij}&= \nabla_{ij}F-F A_{jk}A_{il}g^{kl}+
\nabla_i a^lA_{jl}+\nabla_j a^lA_{il}+a^l\nabla_l A_{ij} \mbox{,}\\
\label{L26}\frac{\partial }{\partial t}H&=\Delta F +F \vert A\vert ^2+g^{ij}a^l\nabla_l A_{ij}.
\end{align}
\end{lemma}

\begin{proof}


Note $\frac{\partial}{\partial t}$ and
$\frac{\partial}{\partial x_i}$  commute. Then,  using \eqref{GW}  and  the fact that $\langle
\frac{\partial f}{\partial x_i}, \nu \rangle =0 $, we obtain
\begin{align}
\frac{\partial}{\partial t}g_{ij}
&=\langle \frac{\partial }{\partial x_i}(F\nu+W),
\frac{\partial f}{\partial x_j}\rangle+\langle \frac{\partial f
}{\partial x_i},
\frac{\partial }{\partial x_j}(F\nu+W)\rangle\nonumber\\
&= F\langle \frac{\partial \nu }{\partial x_i},
\frac{\partial f}{\partial x_j}\rangle + F \langle  \frac{\partial f }{\partial x_i},
\frac{\partial \nu}{\partial x_j}\rangle +\langle \frac{\partial W}{\partial x_i},
\frac{\partial f}{\partial x_j}\rangle+\langle \frac{\partial
f}{\partial x_i}, \frac{\partial W}{\partial x_j}\rangle\nonumber\\
&=-2F A_{ij}+\frac{\partial a^l}{\partial
x_i}g_{jl}+\frac{\partial a^l}{\partial
x_j}g_{il}+a^l\Gamma_{il}^mg_{mj}+a^l\Gamma_{jl}^mg_{mi}\nonumber\\
&=-2F A_{ij}+\nabla_i a^l g_{jl}+\nabla_j a^lg_{il}.\nonumber
\end{align}
This gives \eqref{L21}.

Furthermore, we compute $\frac{\partial}{\partial t}g^{ij}$. Since
$g^{ik}g_{jk}=\delta^i_j$, it follows that
\begin{align}
0=\frac{\partial (g^{ik}g_{jk})}{\partial t}=\frac{\partial
g^{ik}}{\partial t}g_{jk}+g^{ik}\big(-2F A_{jk}+\nabla_j a^l
g_{kl}+\nabla_k a^lg_{jl} \big),\nonumber
\end{align}
and hence
\begin{align}
g^{jm}g_{jk}\frac{\partial g^{ik}}{\partial t}=\frac{\partial
g^{im}}{\partial t}&=-g^{jm}g^{ik}\big(-2F A_{jk}+\nabla_j a^l
g_{kl}+\nabla_k a^lg_{jl} \big)\nonumber\\
&=2F A^{im}-\nabla_j a^i g^{jm}-\nabla_k a^m g^{ik}. \nonumber
\end{align}
Renumbering the indices in the above equation gives \eqref{L22}.

Noting again that $\left <\nu, \frac{\partial f}{\partial x_i}\right >=0$, it  follows  from (2.8) and \eqref{GW}   that
\begin{align}
\frac{\partial}{\partial t} \nu &= \langle \frac{\partial}{\partial
t}\nu, \nu \rangle \nu
+\langle \frac{\partial}{\partial
t}\nu, \frac{\partial f}{\partial x_i}\rangle \frac{\partial
f}{\partial x_j}g^{ij}=-\langle \nu, \frac{\partial}{\partial
t}\frac{\partial f}{\partial x_i}\rangle \frac{\partial f}{\partial
x_j}g^{ij}\nonumber\\
&=-\langle \nu, \frac{\partial}{\partial x_i}(F\nu+W)\rangle
\frac{\partial f}{\partial x_j}g^{ij}
=-g^{ij}\frac{\partial F}{\partial x_i} \frac{\partial
f}{\partial x_j}- a^l g^{ij}A_{il}\frac{\partial f}{\partial x_j}\nonumber\\
&=-\nabla F-a^l g^{ij}A_{il}\frac{\partial f}{\partial x_j}.\nonumber
\end{align}
We consider the evolution of the volume element $d\mu$:
\begin{align}\label{dm}
\frac{\partial}{\partial t} d\mu&=\frac{\partial}{\partial t}
\big(\sqrt{\det g}~ dx\big)=\frac{1}{2}{\sqrt{\det g}}~
g^{ij}\frac{\partial}{\partial t}g_{ij}
~dx=\frac{1}{2}g^{ij}\frac{\partial}{\partial t}g_{ij}
~d\mu  \\
&=\frac{1}{2}g^{ij}(-2F A_{ij}+\nabla_i a^l g_{jl}+\nabla_j
a^lg_{il}) ~d\mu=(-HF+\nabla_l a^l)~d\mu.\nonumber
\end{align}


By using the Gauss  and Codazzi formulas \eqref{GW}, we calculate the evolution of the second fundamental form to obtain \eqref{L24} by
\begin{align}
&\quad \frac{\partial }{\partial t} A_{ij}=\frac{\partial }{\partial
t}\langle \frac{\partial^2
f}{\partial x_i \partial x_j}, \nu\rangle\nonumber\\
&=\langle \frac{\partial^2 }{\partial x_i \partial
x_j}(F\nu+W), \nu\rangle-\langle \frac{\partial^2 f}{\partial
x_i \partial x_j}, \nabla F+a^l g^{mn}A_{ml}\frac{\partial f}{\partial x_n}\rangle \nonumber\\
&=\frac{\partial^2 F}{\partial x_i \partial x_j}+F\langle
\frac{\partial^2 }{\partial x_i \partial x_j}\nu, \nu\rangle+\langle
\frac{\partial^2 }{\partial x_i \partial x_j}W,
\nu\rangle-\Gamma_{ij}^k\langle \frac{\partial f}{\partial x_k},
\nabla F+a^l g^{mn}A_{ml}\frac{\partial f}{\partial
x_n}\rangle\nonumber\\
&=\frac{\partial^2 F}{\partial x_i \partial
x_j}-F A_{jk}A_{il}g^{kl}-\Gamma_{ij}^k\frac{\partial F}{\partial x^k}-\Gamma_{ij}^k
a^lA_{kl}+ \langle \frac{\partial^2
}{\partial x_i \partial x_j}W, \nu\rangle\nonumber \\
&=\nabla_{ij}F-F A_{jk}A_{il}g^{kl}-\Gamma_{ij}^k
a^lA_{kl}+\frac{\partial a^l}{\partial x_j}A_{il}+\frac{\partial
a^l}{\partial x_i}A_{jl}+a^l\frac{\partial
 A_{jl}}{\partial x_i}+ a^l\Gamma_{lj}^k A_{ik}\nonumber\\
&=\nabla_{ij}F-F A_{jk}A_{il}g^{kl}+\nabla_i
a^lA_{jl}+\nabla_j a^lA_{il}+a^l\nabla_i A_{jl}, \nonumber
\end{align}
 where we used the fact that
\begin{align}
\langle \frac{\partial^2 }{\partial x_i \partial x_j}W,
\nu\rangle
&=\langle\frac{\partial }{\partial x_i}\Big(\frac{\partial
a^l}{\partial x_j}\frac{\partial f}{\partial
x_l}+a^l\Gamma_{lj}^k\frac{\partial f}{\partial
x_k}+a^lA_{jl}\nu\Big),\nu \rangle\nonumber\\
&=\frac{\partial a^l}{\partial x_j}A_{il}+ a^l\Gamma_{lj}^k
A_{ik}+\frac{\partial a^l}{\partial x_i}A_{jl}+a^l\frac{\partial
 A_{jl}}{\partial x_i}\nonumber.
\end{align}
Using the above equations, we may compute the evolution of the mean
curvature:
\begin{align}
\frac{\partial }{\partial t}H&=g^{ij}\frac{\partial }{\partial t}A_{ij}+
A_{ij}\frac{\partial }{\partial t}g^{ij}\nonumber\\
&=\Delta F +F \vert A\vert ^2+ g^{ij}\big (\nabla_i a^lA_{jl}+\nabla_j a^lA_{il}+a^l\nabla_l A_{ij} \big)\nonumber\\
&\quad+A_{ij}\big(-\nabla_l a^i g^{jl}-\nabla_l a^jg^{il}\big)\nonumber\\
&=\Delta F +F \vert A\vert ^2+g^{ij}a^l\nabla_l
A_{ij}.\nonumber
\end{align}
These complete the proof of Lemma \ref{L1}.
\end{proof}

\end{document}
